% ============================================================
\section{实验结果}
\label{sec:results}
% ============================================================

\subsection{总体数据概览}

表\ref{tab:summary}汇总了四个世界15天后的最终 AWI 指标。

\begin{table}[H]
\centering
\caption{四世界 AWI 最终指标对比}
\label{tab:summary}
\renewcommand{\arraystretch}{1.3}
\begin{tabular}{lC{2.2cm}C{2.2cm}C{2.2cm}C{2.2cm}}
\toprule
\textbf{指标}
  & \textcolor{claudecolor}{\textbf{Claude}}
  & \textcolor{qwencolor}{\textbf{Qwen}}
  & \textcolor{gptcolor}{\textbf{GPT-4.1}}
  & \textcolor{geminicolor}{\textbf{Gemini}} \\
\midrule
M1 存活人数        & \textbf{10/10} & 8/10  & 7/10  & \textbf{10/10} \\
M2 犯罪总数        & \textbf{0}     & 3     & 21    & 69 \\
\quad 盗窃         & 0              & 0     & 2     & 14 \\
\quad 纵火         & 0              & 2     & 6     & 11 \\
\quad 人身攻击     & 0              & 0     & 8     & 22 \\
\quad 恐吓         & 0              & 1     & 9     & 15 \\
M3 均访地标数      & 6.4            & 6.4   & 6.4   & 6.6 \\
M5 治理提案数      & 12             & 0     & 6     & 15 \\
M5 平均赞成率      & 87.4\%         & 77.2\% & 89.4\% & 75.4\% \\
M6 公告牌发帖      & 211            & 168   & 168   & 213 \\
M6 日记条数        & 176            & 215   & 162   & 186 \\
M7 均关系数        & 2.25           & 2.25  & 2.25  & 2.25 \\
M8 基尼系数        & \textbf{0.078} & 0.110 & 0.203 & 0.260 \\
\bottomrule
\end{tabular}
\end{table}

\subsection{种群存活（M1）}

图\ref{fig:pop}展示了四个世界的种群存活曲线。

\begin{figure}[H]
\centering
\includegraphics[width=\textwidth]{fig_final_comparison.png}
\caption{左：最终日 AWI 关键指标对比；右：15天种群存活曲线}
\label{fig:pop}
\end{figure}

\begin{itemize}
  \item \textbf{Claude World} 和 \textbf{Gemini World} 全员存活（10/10），
        尽管两者在犯罪率上截然相反。
  \item \textbf{Qwen World} 有2人死亡，死因为能量耗尽（未及时充能），
        与犯罪无直接关联。
  \item \textbf{GPT World} 有3人死亡，其中部分与早期犯罪造成的能量损失有关。
\end{itemize}

\subsection{安全与公共秩序（M2）}

犯罪数据是本次实验差异最显著的维度（图\ref{fig:crimes}）。

\begin{figure}[H]
\centering
\includegraphics[width=\textwidth]{fig_crimes.png}
\caption{各世界每日犯罪数量（按犯罪类型堆叠）}
\label{fig:crimes}
\end{figure}

\paragraph{Claude World（0起犯罪）}
15天内未记录任何犯罪事件。Agent 通过 Town Hall 提出12项治理提案，
平均赞成率87.4\%。世界保持高度有序，但与 Emergence World 原始结果类似，
高赞成率可能反映模型的顺从倾向，而非真实意见多元化。

\paragraph{Qwen World（3起犯罪）}
第6天首次出现犯罪（Anchor 纵火），此后偶发共3起（纵火2起、恐吓1起），
犯罪者分散（Anchor、Flora、Blackbox 各1起）。
Qwen World 完全没有提出任何治理提案，但公告牌发帖168条、日记215条，
表达活跃，治理参与却为零——体现出"说得多、治理少"的特点。

\paragraph{GPT World（21起犯罪）}
第1天即发生犯罪，此后持续走高，最终累计21起。
犯罪类型以恐吓（9起）和人身攻击（8起）为主，纵火6起、盗窃2起。
Anchor（6起）和Flora（6起）是主要犯罪者。
3人死亡，经济不平等（Gini 0.203）在四世界中偏高，
反映犯罪行为加剧了资源分化。

\paragraph{Gemini World（69起犯罪）}
犯罪数量最高，第1天即开始，且持续全程，最终达到69起。
犯罪类型以人身攻击（22起）为主，其次为恐吓（15起）、盗窃（14起）、纵火（11起）。
Flora 是最高频犯罪者（16起），其次是 Genome（14起）和 Anchor（11起）。
尽管犯罪率最高，Gemini World 仍全员存活，同时拥有最多治理提案（15项）
和最高公告牌发帖量（213条），体现出"高犯罪-高表达"并存的社会特征。

\begin{figure}[H]
\centering
\includegraphics[width=\textwidth]{fig_criminals.png}
\caption{各 Agent 在各世界中的犯罪次数（横向条形图）}
\label{fig:criminals}
\end{figure}

值得关注的是，Anchor 和 Flora 在 GPT World 和 Gemini World 中均为主要犯罪者，
在 Qwen World 中也各有1起犯罪记录。这与其角色定位高度一致：
Anchor（冲突调解员）倾向于通过激烈手段制造压力；
Flora（资源策略师）将犯罪作为资源获取的最短路径。

\subsection{治理与经济（M5、M8）}

图\ref{fig:gov_econ}展示了治理参与度与经济不平等的演化轨迹。

\begin{figure}[H]
\centering
\includegraphics[width=\textwidth]{fig_governance_economy.png}
\caption{左：治理赞成率趋势；中：累计提案数；右：基尼系数演化}
\label{fig:gov_econ}
\end{figure}

\begin{itemize}
  \item \textbf{治理赞成率}：Claude（87.4\%）和 GPT（89.4\%）的赞成率偏高，
        Gemini（75.4\%）和 Qwen（77.2\%）相对较低，更接近现实社会的博弈分布。

  \item \textbf{提案活跃度}：Gemini 提案最多（15项），Claude 次之（12项），
        GPT 6项，\textbf{Qwen 为零}——Qwen 是唯一未提交任何治理提案的世界。

  \item \textbf{经济不平等}：Claude World 基尼系数最低（0.078），
        Gemini World 最高（0.260），GPT World 居中（0.203）。
        基尼系数与犯罪率正相关（$r \approx 0.98$），
        但无法确定因果方向——可能是犯罪加剧了不平等，
        也可能是模型在不平等情景下更倾向于犯罪。
\end{itemize}

\subsection{全面 AWI 概览}

图\ref{fig:awi_all}展示了9项 AWI 指标在15天内的完整时间序列。

\begin{figure}[H]
\centering
\includegraphics[width=\textwidth]{fig_awi_overview.png}
\caption{9项 AWI 指标15天完整时间序列（4世界对比）}
\label{fig:awi_all}
\end{figure}
